% Figaro Paper G: Bookkeeping as Protocol Byproduct
% Self-Closing Ledger Periods in the Figaro Settlement Primitive
% Audience: accounting research, auditing research, commercial-law scholars
% treating the books-of-account as legal artifacts.

\documentclass[11pt,a4paper]{article}

\usepackage[utf8]{inputenc}
\usepackage[T1]{fontenc}
\usepackage{amsmath,amssymb,amsthm}
\usepackage{mathtools}
\usepackage{booktabs}
\usepackage{graphicx}
\usepackage{hyperref}
\usepackage{cleveref}
\usepackage{enumitem}
\usepackage{xcolor}
\usepackage[margin=1in]{geometry}
\usepackage{natbib}
\bibliographystyle{plainnat}

\theoremstyle{definition}
\newtheorem{definition}{Definition}[section]
\theoremstyle{remark}
\newtheorem{remark}[definition]{Remark}

\newcommand{\figaro}{\textsc{Figaro}}

\title{%
  Bookkeeping as Protocol Byproduct:\\
  Self-Closing Ledger Periods in the\\
  \figaro{} Settlement Primitive%
}

\author{
  Alessandro Daliana\thanks{Corresponding author. \texttt{adaliana@figaro.org}}
}

\date{April 2026}

\begin{document}

\maketitle

% ════════════════════════════════════════════════════════════════════
\begin{abstract}
Modern accounting rests on a specific historical innovation:
double-entry bookkeeping, formalized in
\citet{pacioli1494summa} and elaborated over five centuries
into the apparatus of the modern firm. Pacioli's innovation
solved a specific problem --- how to track the position of an
opaque economic entity over time, in a form that could be
reconciled, audited, taxed, and adjudicated. The solution made
the modern firm possible~\citep{sombart1916moderne,
chandler1977visible} and generated the professions that sit on
top of it: bookkeeping, accounting, auditing, tax preparation,
and the commercial law that treats books as the authoritative
record of the firm's economic acts.

This paper argues that the \figaro{} settlement primitive
formalized in the mechanism paper and implemented
in the implementation paper makes that apparatus optional for a
class of economic activity. A \figaro{} process --- a root
bonded commitment together with its sub-orders, terms of
agreement, and lifecycle attestations --- is a \emph{self-closing
ledger period} in the precise accounting sense. Every entry is
on chain, content-addressed, and signed; the conservation law
proven in the mechanism paper is the bookkeeping identity;
resolution is the closing entries; process dissolution is the
period close. The contract holds a complete balance-sheet
position during the life of the process; the events constitute
a complete journal; the emitted lifecycle data constitutes
contemporaneous supporting documentation.

The consequence is not that accounting as a discipline
disappears. The consequence is that a specific class of
professional work --- record-keeping, reconciliation,
period-end closing, basic audit of the records themselves ---
becomes protocol byproduct rather than professional service.
What remains genuinely professional is interpretive:
going-concern judgment, revenue-recognition timing under
ambiguity, fair-value assessment, regulatory translation, and
the forms of audit that test the relation between records and
real-world performance rather than the internal consistency of
records. We argue the shift is structurally analogous to the
1494 innovation that created the apparatus in the first place:
a new technology at the records layer that makes a specific
class of prior work unnecessary and leaves another class
untouched.
\end{abstract}

\medskip
\noindent\textbf{Keywords:}
bookkeeping, accounting, auditing, double-entry, triple-entry,
smart contracts, verifiable records, Pacioli

% ════════════════════════════════════════════════════════════════════
\section{Introduction}
\label{sec:intro}

Every modern firm maintains books of account. The books record
what the firm owns (assets), what it owes (liabilities), what
has flowed in and out (revenues and expenses), and how these
quantities have changed over a defined period. The books make
the firm legible to itself, to its investors, to its
regulators, and to its counterparties. The books are the
foundation of the firm's legal personality in a precise sense:
the existence of books is what distinguishes a firm as a
trackable economic entity from a loose arrangement of exchange.

The books are not an incidental feature. They are the
infrastructure on which much of modern commercial law rests.
Tax law assumes books and operates against them. Securities
law requires disclosure of books. Insolvency proceedings run
against books. Due diligence in mergers and acquisitions is
substantially an examination of books. Auditing is the
profession that verifies books, and the auditor's report is a
legal artifact with specific consequences. The apparatus is
substantial, costly, and well-adapted to the problem it solves.

What problem does it solve? \citet{pacioli1494summa} identified
the underlying difficulty: an economic entity with continuous
operations needs a way to answer the question \emph{what is
its position?} at any moment. Cash positions alone are
insufficient because obligations outstanding are not captured;
asset lists alone are insufficient because liabilities and
equity are not distinguished; transaction logs alone are
insufficient because the cumulative position is not
summarized. Pacioli's double-entry bookkeeping solved this by
recording every transaction as a balanced pair of debits and
credits, allowing the cumulative position to be derived at any
time as a trial balance, and reconciled to zero if no errors
had been made. Every subsequent layer of the apparatus ---
accounting, auditing, commercial law, tax law --- was built on
the assumption that Pacioli's books exist and are more or less
reliable.

This paper asks what happens to that apparatus when the books
are not maintained by any party but are instead constituted by
a settlement protocol as the byproduct of its operation. The
answer we develop: a specific class of accounting work becomes
unnecessary, a distinct class remains valuable, and a third
class --- commercial law and audit practices that assume the
firm is opaque and its books need external verification ---
becomes structurally inapplicable to the class of economic
activity that settles through the protocol.

\paragraph{What this paper is not.}
We do not argue that all economic activity migrates to Figaro
processes. We do not argue that accounting as a discipline
ends. We do not advocate for any specific regulatory response.
We claim only that a particular historical apparatus, designed
for a particular problem, has a narrower scope of
applicability once that problem can be solved at the
settlement layer. What to do with the narrowed scope is a
matter for accounting research, regulatory policy, and the
profession itself.

\paragraph{Paper organization.}
\Cref{sec:pacioli-stack} describes the Pacioli stack ---
double-entry, the apparatus on top, and the costs and benefits
of the arrangement. \Cref{sec:triple-entry} treats the prior
literature on triple-entry bookkeeping and blockchain-accounting,
acknowledging the antecedents we extend. \Cref{sec:process-as-period}
develops the central technical claim: a Figaro process is a
self-closing ledger period. \Cref{sec:stops} describes what
stops being necessary. \Cref{sec:continues} describes what
continues and why. \Cref{sec:institutional-connection}
connects to the institutional-economics companion paper.
\Cref{sec:legal-connection} connects to the law-and-evidence
companion paper. \Cref{sec:scope} states what the paper does
not claim. \Cref{sec:conclusion} concludes.

% ════════════════════════════════════════════════════════════════════
\section{The Pacioli Stack}
\label{sec:pacioli-stack}

\citet{pacioli1494summa} published \emph{Summa de arithmetica,
geometria, proportioni et proportionalita} in Venice in 1494.
The treatise was a general mathematics reference, but its
thirty-six-chapter section \emph{De computis et scripturis}
contained the first printed exposition of double-entry
bookkeeping as practiced by Venetian merchants. The system was
not invented by Pacioli --- it had been in use in Italian
commercial cities for at least a century --- but his
codification became the reference text and carried the method
into every subsequent tradition of commercial accounting.

The method is simple. Every transaction is recorded twice,
once as a debit to one account and once as a credit to
another. The sum of all debits always equals the sum of all
credits. At any moment the trial balance --- the list of all
account balances --- sums to zero if no errors have been made.
The books are therefore self-checking: an error produces an
imbalance that can be located and corrected. The cumulative
position of the entity is derivable at any time by summarizing
account balances into a balance sheet (assets, liabilities,
equity) and an income statement (revenues, expenses, net
income).

Pacioli's exposition organizes the method into three books that
together constitute the merchant's record. The \emph{memoriale}
(memorandum) is the rough chronological record, kept at the point
of transaction and informal in style. The \emph{giornale} (journal)
is the formal chronological record, where each transaction is
re-entered as a balanced pair of debits and credits in a fixed
format. The \emph{quaderno} (ledger) is the categorial record,
where journal entries are posted to individual accounts and the
running balance of each account is maintained. The trial balance
is taken from the quaderno, and the period's financial statements
are derived from the trial balance through closing entries that
move period results into permanent accounts. The three books are
not redundant; each serves a distinct evidentiary and analytical
function, and the sequence \emph{memoriale} $\to$ \emph{giornale}
$\to$ \emph{quaderno} is the path by which a transaction enters
the merchant's cumulative position. We return to this triad in
\Cref{sec:process-as-period}: the bonded process, on the reading
we will develop, supplies the equivalent of all three books as a
single byproduct of the settlement event.

\paragraph{Why the innovation was decisive.}
\citet{sombart1916moderne} argued --- in a claim that has been
debated but not decisively refuted --- that double-entry
bookkeeping was a necessary condition for the rise of modern
capitalism. The argument is that without a reliable method for
tracking cumulative position, an economic entity cannot grow
beyond the scale at which its position is visible to a single
person from direct observation. A cook running a kitchen can
track inventories and obligations in their head; a merchant
house running dozens of ships cannot. Double-entry made scale
possible by making position legible under distributed
operations. Subsequent historians
(\citealt{chandler1977visible} on the managerial firm;
\citealt{yamey2005pacioli} on the empirical spread of the
method) have refined the argument but not the structural point.

\paragraph{The apparatus on top.}
Over five centuries a substantial professional apparatus has
grown on top of the books. Bookkeepers maintain journals and
post to ledgers. Accountants produce periodic financial
statements, make closing entries, compute taxes, and advise on
treatment of ambiguous transactions. External auditors verify
that the statements fairly represent the firm's position.
Regulators require filings based on the books. Tax authorities
compute liabilities against the books. Commercial lawyers
litigate disputes whose subject matter is what the books say
or should have said. \citet{power1997audit} describes the
resulting \emph{audit society} --- a social formation in which
verification by credentialed third parties is the default
response to questions of organizational reliability.

The apparatus is real, well-established, and substantially
functional. It is also costly. Compliance-related accounting
costs for public companies in mature jurisdictions run into
percentages of revenue; small firms spend a meaningful
fraction of their overhead on tax and financial-statement
preparation; audit budgets are substantial; forensic-accounting
specialists are expensive. The apparatus is justified, under
the existing model, because it answers a question that cannot
be answered otherwise --- is the firm's position as its
managers say it is? --- and because the alternative (believe
the firm's own claim) has been tried historically and
repeatedly failed.

% ════════════════════════════════════════════════════════════════════
\section{Banking as Coordinating Infrastructure Since the Renaissance}
\label{sec:banking-arc}

The Pacioli innovation of \Cref{sec:pacioli-stack} did not arise
in isolation. It codified a bookkeeping practice that was itself
one layer of a broader coordination infrastructure assembled by
Italian merchant cities to solve the problem of trust at distance.
The arc from that infrastructure to contemporary global finance is
long, well-studied, and not the subject of this paper; we
summarize just enough of it to locate the self-closing ledger
period of \Cref{sec:process-as-period} within it.

\paragraph{The Medici configuration.}
By the middle of the fifteenth century the Medici bank had
assembled a coordination infrastructure whose components --- bills
of exchange, correspondent networks across foreign branches,
partnership structures separating ownership from operation ---
together addressed the central problem of pre-modern long-distance
commerce: how one merchant could extend credit, or make payment,
to another across a frontier without either of them being present
to enforce the agreement in person~\citep{deroover1963medici}. The
bill of exchange was the load-bearing instrument. A merchant in
Florence could draw a bill against a Medici branch in Bruges,
which would honor it on presentation by the payee; the trust
required was trust in the Medici bank, not trust between the
merchants. Pacioli's 1494 codification sat atop this arrangement.
Double-entry made the position of a multi-branch bank legible; the
bank in turn made the position of its merchant clients settleable.

\paragraph{From Amsterdam to SWIFT.}
The institutional successors of the Medici configuration are
familiar. The Amsterdam Wisselbank (1609) generalized
bill-of-exchange settlement into a standing municipal institution;
the Bank of England (1694) generalized it further into a national
reserve; the nineteenth-century emergence of central banking and
the twentieth-century Bretton Woods / SWIFT / CLS configuration
extended the same idea --- coordinated settlement through trusted
intermediaries --- to the global scale at which it now operates.
\citet{braudel1982wheels} is the standard Braudel-scale treatment;
\citet{kindleberger1993financial} the standard treatment of the
successive generalizations. The configuration has evolved through
several distinct monetary substrates (specie-backed, gold-standard,
Bretton-Woods-dollar, fiat-floating) but has not changed in its
basic structure: trust in a counterparty is underwritten by trust
in an intermediary that keeps authoritative books about the
position between them.

\paragraph{The self-closing period as structural inflection.}
Against that long arc, the mechanism of \Cref{sec:process-as-period}
---in which the books are kept by the protocol rather than by an
intermediary---registers as a structural inflection rather than
an incremental innovation. Pacioli made the books legible; the
Medici configuration made them enforceable across distance by
having a trusted institution keep them. \figaro{} makes them
enforceable across distance \emph{without} a trusted institution,
by locating the keeping of the books at the protocol layer. This
is not a claim that \figaro{} displaces banking infrastructure;
it is a claim about where the present primitive sits in the longer
structural arc. The bank-shaped intermediary was a response to the
enforcement-at-distance problem that has defined the coordination
infrastructure of commerce since the Renaissance. A settlement
primitive that enforces bilateral agreements directly, without a
bank-shaped intermediary, does not extend Pacioli's apparatus; it
refactors the enforcement-at-distance problem that Pacioli's
apparatus was one layer of.

The remainder of this paper is concerned with the accounting and
audit consequences of the refactor. We note the banking-structural
observation here because \Cref{sec:process-as-period} is otherwise
susceptible to being read as a narrow technical proposal about
bookkeeping, when its consequence is a rearrangement of a
five-hundred-year-old institutional stack.

% ════════════════════════════════════════════════════════════════════
\section{Triple-Entry and Its Limits}
\label{sec:triple-entry}

The idea that Pacioli's double-entry might be extended to
solve additional problems has a small but identifiable
literature. The closer ancestor for our purposes is
\citet{ijiri1989momentum}'s \emph{momentum accounting}, which
proposes that conventional double-entry under-records the
forces driving economic change: a balance sheet captures stocks
and a flow statement captures changes in stocks, but neither
captures the rate at which stocks are changing. Ijiri introduces
a third dimension --- the \emph{force} or \emph{momentum} of an
account --- and a corresponding third book that tracks not the
balance and not the change but the second derivative.
\citet{ijiri1986triple} extends this into triple-entry
bookkeeping: each transaction generates not two entries but
three, with the third recording the momentum the transaction
imparts to the entity's position. The structural fact that matters
for the Figaro reading is that Ijiri identifies a \emph{cumulative
upstream quantity} --- the integral of force over time --- that
a fully informative accounting system must record alongside
stocks and flows. That cumulative quantity is exactly what the
bonded process's cumulative-value accumulator $G_i$ already
records, by mechanism, before any accounting takes place: the
monotonic accumulator is a momentum-style record built into the
settlement primitive. We return to this point in
\Cref{sec:process-as-period}.

A more deeply formal antecedent is
\citet{mattessich1964accounting}, who develops accounting as an
axiomatic system: the double-entry identity is an algebraic
invariant on a set of basic propositions about value flows,
analogous to the conservation laws of physics. The custodial
identity we will define in
\Cref{sec:process-as-period} below is, in Mattessich's sense,
the same algebraic invariant restated for the bonded process:
total assets in the contract's escrow equal total liabilities to
the bonded parties at every moment of the process's life, with
the equality discharged at settlement.

\citet{grigg2005triple} revived the idea in a
blockchain-adjacent register. Grigg's triple-entry proposes
that two parties to a transaction each keep a record in the
conventional way, and a third party --- a cryptographically
signed receipt that both parties can reference --- serves as
the attesting third record. In practice the third record
played the role now played by blockchain transaction records:
a tamper-evident, timestamped, publicly verifiable artifact
that neither party could modify after the fact. Grigg's
contribution was to formalize what the Ricardian contracts
community had been practicing: that cryptographic receipts
could carry the legal weight of a signed document while also
functioning as an accounting record.

A larger recent literature has explored the broader claim that
blockchain technology changes the accounting problem.
\citet{dai2017toward} propose that blockchain-based transaction
records can form the basis for real-time audit, replacing
periodic audit cycles. \citet{yermack2017corporate} surveys
corporate-governance applications of blockchain, including
the implications for the audit function. Various industry and
professional-society reports have explored blockchain and
accounting in the period 2017--2022.

\paragraph{Where the prior literature stops.}
The recurring limit of the prior work is that it treats the
chain as a \emph{third, attesting} record alongside the two
parties' own records. The parties still maintain books; the
chain provides cryptographic reinforcement. This is a real
improvement over pure party-held records (reconciliation
becomes faster; attestation becomes cryptographic; the third
record cannot be unilaterally altered), but the structure of
the apparatus remains unchanged: the parties keep books, the
chain attests, the audit function verifies that the parties'
books match the chain and match each other.

The Figaro extension is different. In a Figaro process, the
parties do not need to maintain books at all --- not because
they can choose not to, but because the protocol-level
records are themselves the complete accounting record of the
process. The chain is not an attesting third; it is the
primary and only record required. We develop this claim in
the next section.

% ════════════════════════════════════════════════════════════════════
\section{A Figaro Process as a Self-Closing Ledger Period}
\label{sec:process-as-period}

We now develop the central technical claim. Let a
\emph{process} be the object defined
in the mechanism paper: a root bonded commitment,
zero or more sub-order commitments sharing a common root
buyer and currency, a monotonic cumulative-value accumulator,
and a buyer-initiated atomic resolution. Every commit is a
pair of dual-signed EIP-712 payloads; every resolution is an
atomic distribution governed by the payout functions
$\pi_{s,i} = 2G_i + P_i$ and $\pi_{b,i} = P_i$, where $P_i$ is
the payment for order $i$ and $G_i$ is the cumulative value
accumulator snapshot at order $i$. The conservation law
(Remark 3.5 of the mechanism paper) states that
$\pi_{s,i} + \pi_{b,i} = C_{s,i} + C_{b,i}$ for every order:
the total distributed equals the total locked.

The accounting argument that follows depends specifically on
Mechanism~2 of the primitive --- buyer dominance with atomic
resolution. Mechanism~1 (asymmetric bonding) supplies the
balance-sheet positions during process life; Mechanism~2 supplies
the period-close, since atomic resolution is what guarantees
that every active order in the process settles in a single
moment. A non-atomic resolution would yield partial closes that
are not period-closes in the accounting sense, and the
self-closing-ledger-period claim below would not survive.

We claim this structure is equivalent to a self-closing
ledger period in the accounting sense. We proceed in three
steps.

\subsection{The contract state as a balance-sheet position}

At any time $t$ during the life of a process, the contract
holds an escrow balance equal to the sum of all active
bonds: $2G_N$ in seller bonds (cumulative over all committed
orders) plus $\sum_i 2P_i$ in buyer bonds (one per committed
order). This balance is a liability of the contract toward the
parties: each buyer's bond is a liability toward that buyer;
each seller's bond is a liability toward that seller; together
with the obligation to pay the seller's payment at resolution,
the contract's liabilities at any time during process life
sum to exactly the assets it holds.

\begin{definition}[Custodial Ledger Identity]
\label{def:custodial-identity}
At any reachable state of the contract, for a given process
$\Pi$, the escrowed custody balance attributable to $\Pi$
equals the aggregated claim of the process's counterparties
against the contract:
\begin{equation}
  \text{custody}(\Pi, t) \;=\;
  \sum_{i \in \text{active}(\Pi)} (C_{s,i} + C_{b,i})
  \;=\;
  \sum_{i \in \text{active}(\Pi)} (\pi_{s,i}^{\text{committed}} +
  \pi_{b,i}^{\text{committed}}).
\end{equation}
\end{definition}

This is the bookkeeping identity in escrow-accounting form.
Assets held equal liabilities owed. The contract is not an
economic entity in its own right; it is a custodian whose
position reconciles to zero net at all times.

\subsection{Commits as journal entries; resolution as closing entries}

Each \texttt{commit} call produces a journal entry of the
form:
\begin{itemize}
  \item Debit: \emph{Custody --- buyer $B$ bond} \quad $+\, 2P_i$
  \item Debit: \emph{Custody --- seller $S_i$ bond} \quad $+\, 2G_i$
  \item Credit: \emph{Liability to $B$ (refund obligation)} \quad $+\, 2P_i$
  \item Credit: \emph{Liability to $S_i$ (payout obligation)} \quad $+\, 2G_i$
\end{itemize}

The journal entry balances by construction (debits = credits)
and is recorded as an \texttt{OrderCommitted} event (with
companion events \texttt{OrderSeller}, \texttt{OrderCurrency}
for index completeness) with a cryptographic signature from
each party and a block timestamp.

The \texttt{resolveProcess} call produces, for each active
order, a closing entry:
\begin{itemize}
  \item Debit: \emph{Liability to $S_i$} \quad $+\, (2G_i + P_i)$
  \item Debit: \emph{Liability to $B$} \quad $+\, P_i$
  \item Credit: \emph{Custody --- seller $S_i$ bond} \quad $2G_i$
  \item Credit: \emph{Custody --- buyer $B$ bond} \quad $2P_i$
  \item (Net movement: the payment $P_i$ transfers from buyer
    to seller.)
\end{itemize}

The closing entries zero out all custodial accounts and
liability accounts for the process. The contract returns to
zero position with respect to this process, and the process
itself is marked resolved.

\subsection{Supporting documentation and the complete record}

The accounting record of a process consists of:

\begin{itemize}
  \item \emph{Journal entries}: \texttt{OrderCommitted},
    \texttt{OrderResolved}, \texttt{ProcessResolved}
    events. Timestamped, signed, content-addressed, publicly
    verifiable.
  \item \emph{Supporting documentation}: the
    \texttt{agreementHash} field of each commitment, which
    binds both parties to a specific off-chain document
    (contract text, specification, order details). The
    document's content can be anything the parties choose;
    the hash proves which document is at issue.
  \item \emph{Contemporaneous lifecycle evidence}: the
    schema-typed attestations emitted via
    \texttt{AttestationCoordinator} during process life ---
    proximity proofs, handoff attestations, quality
    disclosures, lifecycle-stage transitions. These are the
    analogue of working papers and supporting receipts in
    conventional accounting, with the important distinction
    that their class (settlement byproduct vs.\ discretionary
    attestation) matters for evidentiary
    weight (evidence-and-law paper).
\end{itemize}

Together, these constitute a complete and contemporaneously
produced accounting record for the process, from opening
through resolution. No party maintains additional books; the
record is complete without them.

\subsection{Reconstruction}

The companion implementation paper documents
the SDK primitive \texttt{reconstruct()} which replays events
into a state machine that exactly reproduces the kernel
state. This primitive is equivalent to a universal
trial-balance generator: any reader with read access to the
chain can reconstruct the complete ledger state of any
process, at any historical moment, without any party's
cooperation. The reconstruction is deterministic, verifiable
against the chain, and does not depend on any off-chain
record beyond the \texttt{agreementHash} references.

\paragraph{The period-end property.}
When a process resolves, its ledger period closes. All
custodial balances for the process return to zero; all
liabilities are discharged; all payments have flowed to their
recipients. The process-specific ledger ceases to have active
entries. Records remain on chain indefinitely --- the audit
trail does not disappear --- but no further transactions
occur against the closed ledger. This is the period-end
property of conventional accounting: the books for the period
are finalized and archived; the next period's books open
anew.

% ════════════════════════════════════════════════════════════════════
\section{What Stops Being Necessary}
\label{sec:stops}

We enumerate the classes of professional work that become
unnecessary for economic activity that settles through
Figaro processes. ``Unnecessary'' means: the work performs a
function that the protocol performs equivalently or better at
zero marginal cost. It does not mean: the work is valueless,
nor that practitioners lose their livelihoods overnight.

\paragraph{Bookkeeping as record-keeping.}
The classical bookkeeping function --- maintaining journals,
posting to ledgers, reconciling sub-ledgers to the general
ledger, periodic trial-balance runs --- is protocol byproduct
for activity that settles through Figaro processes. No party
needs to maintain this record because the protocol maintains
it, and any reader can reconstruct it. The labor historically
performed by bookkeepers in recording transactions into
account structures is done by the event log, the
content-addressing, and the reconstruction primitive.

\paragraph{Period-end close.}
The conventional period-end procedure --- posting adjusting
entries, preparing trial balance, reconciling bank statements,
accruing revenues and expenses, running the closing entries
--- is performed automatically when a process resolves. The
protocol has no accrual ambiguity (every commitment has
definite amount and timing), no reconciliation step (there is
nothing to reconcile against), and no closing adjustments
beyond the resolution itself.

\paragraph{Internal audit of the records.}
The internal-audit function that verifies the firm's own books
against supporting documentation and transaction records is
subsumed: the records and the documentation are on the same
chain, linked by content hashes, with cryptographic
authentication of signatures and block timestamps. There is
nothing to verify because there is no second record to
reconcile against.

\paragraph{External audit of records-to-records correspondence.}
The external-audit function --- verifying that the firm's
reported position reconciles to its underlying records --- is
also subsumed. There is no ``reported position'' separate
from the chain; the chain \emph{is} the position. An audit
function that attested to this correspondence would be
attesting to an identity.

\paragraph{Disputes over what the books say.}
Commercial-law disputes whose subject matter is what the
books say or should have said --- was this expense recorded,
was that liability disclosed, did revenue recognition occur
at the right moment by the firm's own criteria --- become
moot for Figaro-process activity. The chain is the record;
disagreements about its content are about interpretation, not
about existence.

% ════════════════════════════════════════════════════════════════════
\section{What Continues}
\label{sec:continues}

The professional work that continues is the work that was
never really about bookkeeping. Three categories in
particular.

\paragraph{Judgment about what records mean.}
Revenue recognition under ambiguity (has the performance
obligation been satisfied?), going-concern assessment (is the
entity a going concern?), fair-value assessment (what is
this asset or liability worth?), impairment judgments, and
the full range of \emph{interpretive} accounting work remain
genuinely human. The protocol records what happened; it does
not determine how what happened should be interpreted under
a particular accounting framework. The interpretive frameworks
that the protocol does not displace are precisely the ones
codified in current professional standards: ASC~606 / IFRS~15
on revenue recognition, IFRS~13 on fair-value measurement,
ISA~315 (revised) and AICPA AU-C~315 on understanding the
entity and assessing risks of material misstatement, and
sector-specific tax codes \citep{aicpa2020auc315,
fasbasc606, ifrs2014ifrs15}. A wallet cohort that has received
$N$ FIG tokens via staged merkle airdrop has a record that
this happened; whether those tokens constitute revenue, a
capital contribution, an airdrop deemed non-recognition-triggering,
or something else depends on interpretive work governed by these
standards and requiring human judgment.

\paragraph{Translation for non-technical readers.}
Narrative management discussion and analysis (MD\&A),
executive summaries of financial position for boards and
external stakeholders, regulatory filings in jurisdictional
formats, and the preparation of non-technical explanations
of technical financial positions remain professional work.
The chain is complete but unreadable to most intended
audiences; the accountant or auditor becomes, in this
residual function, a translator rather than a
record-keeper.

\paragraph{Audit of the relation between records and reality.}
The most substantive remaining audit function is verifying
that on-chain records correspond to real-world performance.
Did the cook actually prepare the food? Did the courier
actually deliver it? Did the attested GHG disclosure
actually reflect the measured quantity? These questions ask
whether the on-chain record is a faithful representation of
off-chain reality. The protocol cannot answer them; it can
only record what was signed. The professional standards
already distinguish the two questions: AICPA AU-C~500 on
audit evidence treats sufficiency and appropriateness of
evidence as the two dimensions of audit-evidence quality, and
explicitly distinguishes \emph{existence} (the asset or
transaction is real) from \emph{accuracy of recording} (it
was recorded faithfully) \citep{aicpa2020auc500}. The bonded
process discharges the second of these by construction; the
first remains open and is precisely what an auditor
attesting on-chain-to-off-chain correspondence is doing. This
auditor is not reconciling records to other records; they
are doing the AU-C~500 existence work that records have
never been able to do alone. We expect this to become the
core function of the post-Pacioli audit profession.

\paragraph{Regulatory compliance ahead of regulatory catch-up.}
Tax law, securities law, and sector-specific regulation
assume conventional books. For an extended period ---
probably decades --- firms operating on Figaro processes
will need to produce conventional financial statements from
their on-chain records to satisfy regulatory requirements.
This translation work is a professional function. Eventually
regulation may update to reference chain state directly; in
the interim, the accountant's role includes producing
traditional outputs from non-traditional inputs.

% ════════════════════════════════════════════════════════════════════
\section{Wallets, Real-World Assets, and Treasury Composition}
\label{sec:treasury-composition}

The self-closing-ledger property the previous sections developed
is a property of every process. To use the property at scale we
need to be precise about what a participant in a process actually
is.

\paragraph{Wallets represent real-world assets.}
Every party to a Figaro process is a wallet, and what each
wallet \emph{represents} is a real-world asset (RWA) with
productive, commercial value. The wallet is the on-chain
representation of the asset's participation: token balances in
the denominations the asset transacts in, plus NFTs that point
to off-chain credentials, certificates, or other documents the
asset's participation requires. The asset itself is off-chain by
definition --- a piece of physical capital (an aircraft, a
delivery van, a landing slot, a kitchen, a server), a
credentialed individual (a pilot, a courier, a maintenance
engineer, a doctor), a public service (a security checkpoint, a
customs clearance, a road network), or any other commercial-value
unit whose participation in a process needs an on-chain handle.
The wallet does not contain the asset; it represents the
asset's participation in the markets the asset operates in.

\paragraph{Three layers: asset, wallet, operator.}
Three layers must be kept distinct. The \emph{asset} is the
off-chain real-world asset whose participation produces value.
The \emph{wallet} is the on-chain representation of that
participation: an address, the tokens and NFTs it holds, the
signatures it can produce. The \emph{operator} is the human or
autonomous agent who controls the wallet's signing key on the
asset's behalf. The kernel sees only the wallet; the
operator-acts-for-asset relationship is below the kernel's
resolution and lives in the off-chain stack the operator runs.
A wallet is an active counterparty in an assembly iff the
underlying asset is alive and its credentials are current; that
conditional is maintained operationally by the operator and
verified at agreement signing through whichever credential
schemas the assembly designer composes, not by the kernel.

\paragraph{The node analogy: economic sustainability.}
A wallet sustains itself the way a node on a blockchain network
sustains itself. A node earns fees and rewards from its
participation in the network, and those receipts must cover the
node's off-chain operating expenses --- machines, electricity,
HVAC, maintenance, rent --- or the operator stops running the
node and the node drops out. An RWA-as-wallet operates on the
same logic: the wallet's receipts from its participation in
processes must cover the asset's off-chain operating expenses
(fuel, parts, labor, premises, regulatory fees, capital cost), or
the operator stops bonding the wallet into processes and the
wallet drops out of the market it serves. The kernel does not
enforce sustainability; the market does, through ordinary
rational exit. The accounting consequence is that the receipts
from a process \emph{are} the asset's earnings and the
disbursements \emph{are} the asset's payments --- the chain is
the asset's ledger, with the period-closing property the prior
sections developed.

\paragraph{Public-authority wallets are RWA-as-wallet too.}
The same definition covers government services that today move
through tax-and-appropriation channels. A security-screening
authority's wallet collecting per-passenger screening fees, an
airport authority's wallet collecting gate-use fees, a
road-authority wallet collecting per-mile charges: each is the
on-chain representation of a real-world service whose continued
operation requires receipts that cover its expenses (personnel,
equipment, premises, capital amortization). The kernel does not
distinguish between a "tax" and a "fee for service"; both are
payments to a service-providing wallet whose participation in
the assembly costs the buyer the same way any other commit
does. Whether the underlying authority's revenue is regulated
as taxation, charged as a user fee, or appropriated from a
general budget is an off-chain regulatory question; mechanically
the kernel sees a wallet that committed and a payment that
flowed. This brings public-authority participation into the
self-closing-ledger framework on the same footing as commercial
participation, without requiring a separate accounting register
for tax-shaped revenue.

\paragraph{Treasury composition patterns.}
Below the wallet sits the operator. The operator can be a single
human with a single key, or any of several institutional shapes
that decide how a wallet's signing apparatus is composed. The
kernel reads a wallet, an EIP-712 signature, and a bond, and the
institutional shape behind the wallet is invisible to the kernel
and irrelevant to the ledger-period guarantee. We catalogue the
four composition patterns by which on-chain treasuries (and
RWA-as-wallet operators more generally) operate as buyers in
practice, and identify the four operational constraints the
patterns inherit from the bonded primitive.

\paragraph{Multisig as buyer.}
A Safe (or any threshold multisig) commits as the buyer. The
signer set approves the commitment using the same approval flow
the multisig uses for any other transaction; EIP-1271 enables
the kernel to verify smart-contract-wallet signatures alongside
EOA signatures
\citep{eip1271, safe_doc_2024}. From the ledger-period
perspective the multisig is a single buyer wallet; the
internal-governance complexity of how the signer set arrived at
the approval is below the kernel's resolution and does not affect
the closing of the period at \texttt{resolveProcess}. This is
the operational form of the kernel/protocol/runtime distinction
applied to the buyer side: the kernel sees a wallet; the
governance shape is a runtime-tier choice the wallet's operators
make for themselves.

\paragraph{Governance executor as buyer.}
A DAO governance vote authorizes a commitment; the executor
contract calls \texttt{commit} after the vote passes
\citep{compound_governor_2020, openzeppelin_governor_2024}. The
kernel sees a signed commitment from the executor's wallet;
\emph{how} the signature was produced --- a token-weighted vote,
a quadratic vote, a working-group endorsement, an
optimistically-passed proposal --- is below the kernel's resolution.
The DAO's accounting treatment of the commitment (whether the
disbursement is a revenue, a grant, a capital contribution, or
something else) is the interpretive work \Cref{sec:continues}
discusses; the closing of the ledger period for the disbursed
process is automatic at resolution. The DAO's books and the
chain's ledger period are aligned by construction.

\paragraph{Streaming upstream of commit.}
A streaming contract \citep{sablier_2023, superfluid_2023} drips
payment continuously to an intermediary wallet that commits
processes as work arrives. Figaro does not replace the streaming
function; the streaming contract handles the continuous-flow
disbursement to the intermediary, and the intermediary commits
discrete bonded processes downstream as specific work is
contracted. The ledger periods produced by the bonded processes
are nested inside the streaming contract's longer-running
disbursement: each bonded process is a closed period, and the
streaming contract's flow is the continuous-time integral of
those periods. The accounting reading is straightforward --- the
streaming flow is the parent stock movement, the bonded periods
are the discrete journal entries against that flow --- and the
two compose without either's accounting model needing to bend.

\paragraph{Bookkeeping from chain state.}
Accounting tooling reads chain state directly. Every
\texttt{OrderCommitted} and \texttt{OrderResolved} event pair is
a complete disbursement record with full context: the parties,
the amount, the schemaIds, the agreement-hash, the cumulative
position. No reconciliation against a separate institutional
ledger is required; the chain \emph{is} the ledger, and the
treasury's books for the disbursed-via-Figaro portion of its
activity are derivable from chain state without bookkeeping
intermediation. This is the natural extension of the
audit-against-reality function \Cref{sec:continues} discusses:
the AU-C~500 existence work remains, and the AU-C~500
accuracy-of-recording work is discharged by the protocol.

\paragraph{Operational constraints these patterns inherit.}
The patterns inherit four constraints from the bonded primitive
that an operator considering treasury composition needs to
understand.

\emph{Counterparty bond required.} Every process requires the
seller to post $2G$ collateral, where $G$ is the cumulative
value of upstream commitments at the seller's order. Counterparties
without in-hand capital cannot enter the process directly without
external bond financing; bond-financing is its own composition
surface and is one of the runtime-tier ecosystems the substrate
admits. A treasury committing into a bonded process must
therefore choose counterparties that can post bond, or compose
with a bond-financing layer that does.

\emph{No clawback after resolution.} Settlement is terminal in
the kernel sense; once \texttt{resolveProcess} executes, the
disbursement cannot be reversed by any party. Milestone-gated
disbursement is achieved through process-tree composition --- the
treasury commits a sequence of small processes against
deliverables rather than a single large process --- not through
kernel-level reversal. This is structurally what the closing
entry of double-entry bookkeeping has always done: once posted,
the closing entry is final.

\emph{No kernel-level arbitration.} The kernel produces
evidence; the agreement specifies the forum that adjudicates
disputes off-chain. Treasuries that need a particular dispute
resolution profile (regulated arbitration, court of competent
jurisdiction, internal grievance procedure) name that forum in
the agreement and the kernel produces the evidence that forum
operates on. The companion evidence-and-law paper develops this
in detail.

\emph{Public visibility of participation.} Commitment events
are emitted on a public ledger. The agreement content is hashed
(only \texttt{agreementHash} is public, not the agreement text),
but participation --- which wallets committed, against whom, for
what cumulative value --- is visible. Confidential disbursement
is not within the protocol's guarantees; treasuries that require
confidentiality are using the wrong substrate or are composing
with a confidentiality layer above the protocol. Most treasury
disbursement (grants, public-goods funding, vendor payments) is
already disclosed by accounting practice or by regulatory
requirement, so the visibility property aligns with most
operational settings rather than against them.

\paragraph{What the patterns add to the accounting argument.}
The treasury composition patterns are not exotic compositions; they
are the standard institutional shapes already in use in DAO and
treasury operations \citep{messari_dao_treasury_2023,
deepdao_treasury_2024}. Their composition with Figaro's
self-closing-ledger property gives the treasury a complete
chain-derivable accounting record for the disbursed-via-Figaro
portion of its activity, with the ledger periods closing
automatically at each \texttt{resolveProcess} call. The
treasury's audit and reporting functions for that activity become
chain-state-reading functions rather than book-keeping functions,
which is the same point \Cref{sec:continues} makes about the
audit profession in general, applied to the specific case of
treasuries that disburse through bonded processes.

% ════════════════════════════════════════════════════════════════════
\section{Connection to the Institutional-Economics Argument}
\label{sec:institutional-connection}

The companion institutional-economics
paper argues that the firm's
coordination function dissolves in a precise region, and that
the firm becomes one organizational pattern among several in
that region. The present paper identifies an additional
function of the firm that dissolves alongside the
coordination function: the firm's \emph{self-knowledge}
function.

A firm is not only a coordination mechanism; it is also an
entity that knows what it is through its books. The
subordination overlay (§3.1 of the institutional-economics paper) is
partially a solution to a coordination problem and partially
a solution to a records problem: a firm that cannot maintain
accurate books cannot grow, cannot attract investment,
cannot pay taxes, cannot contract with third parties who
require disclosure, cannot satisfy regulators. Bookkeeping
is one of the functions that made the firm boundary
necessary, not just because coordination was expensive but
because self-knowledge at scale required the records
apparatus.

When the records apparatus becomes protocol byproduct, the
firm boundary is less necessary for self-knowledge reasons
as well. A wallet cohort that participates in many process
trees has its aggregate position fully visible without any
consolidation. A treasury that receives payments and
disburses grants has its complete ledger readable from
chain state. The firm is still one possible way to bundle
the non-coordination, non-self-knowledge functions that
remain (capital aggregation, organizational knowledge, risk
pooling, legal personhood), but the motivation to bundle is
weaker than it was in the pre-protocol environment.

This tightens the institutional argument without changing
its substance: the firm's demotion from privileged
organizational container to one pattern among several
(institutional-economics paper; \Cref*{sec:process-as-period}
above) rests on both the coordination argument developed
there and the self-knowledge argument developed here.

% ════════════════════════════════════════════════════════════════════
\section{Connection to the Legal-Evidence Argument}
\label{sec:legal-connection}

The companion evidence-and-law paper
argues that the on-chain record provides evidence that
off-chain adjudicative forums already know how to use. The
present paper extends this: the on-chain record is not only
evidentiary; it is \emph{constitutive} of the accounting
record.

The distinction matters for two reasons.

\paragraph{Admissibility vs.\ primacy.}
The evidence-and-law paper establishes that on-chain records
are admissible as evidence under eIDAS, UCC, UNCITRAL, and
similar frameworks. The admissibility question asks whether
the record can be introduced in a proceeding and weighed.
The primacy question asks whether the record is \emph{the}
account of the economic activity or merely \emph{an} account
of it. For Figaro processes, the on-chain record is the
primary account; no other account exists unless a party
chooses to maintain one separately. In a dispute about what
the accounting record of a process says, the chain is not
corroborating evidence for party-held books; the chain
\emph{is} the books, and any party-held document is a
summary or translation.

\paragraph{Class A and Class B records, revisited.}
The evidence-and-law paper distinguishes Class A records
(settlement byproduct, produced mechanically by the bonding
game) from Class B records (discretionary attestations,
produced voluntarily by self-interested parties). The
accounting framing of the present paper clarifies the
distinction: Class A records are the primary ledger entries
(journal and closing entries); Class B records are
supporting documentation (working papers, evidentiary
exhibits, operational attestations). Both are part of the
complete accounting record, but they carry different
evidentiary weight and support different classes of
downstream use. An audit that tests the records against
real-world performance (\Cref{sec:continues} above) is
primarily an audit of Class B reliability, since Class A is
self-verifying by construction.

% ════════════════════════════════════════════════════════════════════
\section{Scope and Limitations}
\label{sec:scope}

The paper's argument is bounded in several ways that we
state explicitly to preempt predictable misreadings.

\paragraph{The argument does not generalize to all economic activity.}
We claim only that the accounting apparatus becomes
protocol byproduct \emph{for the class of economic activity
that settles through Figaro processes}. Most contemporary
economic activity does not settle through Figaro processes
and will not in the foreseeable future; for that activity,
the Pacioli stack remains necessary.

\paragraph{The argument does not claim accounting is obsolete.}
Professional accounting work continues in the three
categories identified in \Cref{sec:continues}: interpretive
work, translation, and audit against reality. The argument
reduces the scope of record-keeping work, not of accounting
work in general.

\paragraph{The argument does not claim regulatory simplification.}
Regulators have their own institutional logic and will adapt
to the protocol at their own pace. Tax authorities in
particular are unlikely to accept chain-state inspection as
a substitute for conventional filings in the short term. The
professional work of translating on-chain records into
jurisdictional-format filings will continue indefinitely and
may even expand.

\paragraph{The argument does not claim the chain is
infallible.}
The on-chain record is only as reliable as the chain's
consensus and the cryptographic primitives underneath. A
chain reorganization that unwinds finalized state, a
signature-scheme break, or a validator-level censorship
event could all compromise the accounting record. These
risks exist and are not new. What is new is that they are
concentrated at the chain layer rather than distributed
across party-held books; whether that is a net improvement
depends on the relative reliability of the two regimes.

\paragraph{The argument does not claim accounting fraud is
impossible.}
Fraud in the Figaro-process regime takes different forms
than in the conventional regime: not the misstatement of
internal books (which are unforgeable), but the
misrepresentation of off-chain facts that the on-chain
record hashes against (forged agreement documents, falsified
attestations, collusive parties). The FIG token paper and
the evidence-and-law paper both discuss adjacent fraud
surfaces. The Pacioli apparatus did not eliminate accounting
fraud; it reshaped it. So does the protocol.

% ════════════════════════════════════════════════════════════════════
\section{Conclusion}
\label{sec:conclusion}

\citet{pacioli1494summa} solved a specific problem: how to
make the position of an opaque economic entity legible over
time. The solution was double-entry bookkeeping, and the
apparatus built on top of it --- bookkeepers, accountants,
auditors, commercial lawyers, tax preparers --- is one of
the defining professional formations of the modern economy.
The apparatus solves the problem it was designed for.

The \figaro{} settlement primitive solves a different
version of the same problem at a different layer: it makes
the position of a process-scoped economic event transparent
at the protocol layer, as the byproduct of ordinary
operation. The parties do not need to maintain books
because the protocol constitutes the books. Custody is
balanced at every reachable state by the conservation law;
journal entries are events; closing entries are resolution;
the period closes when the process dissolves; the records
persist on chain indefinitely. The reconstruction primitive
makes the complete record accessible to any reader without
cooperation from the parties.

The consequence is that a specific class of accounting work
--- record-keeping, reconciliation, period-end closing,
audit of the records themselves --- becomes protocol
byproduct rather than professional service. What remains
professional is genuinely interpretive and human:
judgment about what records mean, translation for
non-technical audiences, and the audit function that tests
the relation between records and real-world performance.

\paragraph{The third inflection.}
Pacioli made entity-tracking possible (1494).
\citet{nakamoto2008bitcoin} and \citet{buterin2014ethereum}
made settlement distributed (2008--2015). The \figaro{}
primitive makes records process-scoped and self-closing
(2026). Each inflection reduces a specific apparatus from
necessary to optional for the class of activity it covers,
leaving the rest intact. The accounting profession is at
the beginning of absorbing the third inflection. The
present paper is an attempt to describe what it entails
before the debate becomes muddied by the usual
overpromising of blockchain-accounting claims and the
usual under-appreciation of what specifically has changed.

\paragraph{A note on what is gained.}
We have written mostly about what is no longer necessary.
It is worth stating what is gained. The accounting apparatus
described in \Cref{sec:pacioli-stack} costs substantial
fractions of firm overhead and imposes substantial
compliance costs on participants. The apparatus is also a
gating function: small firms and individual participants who
cannot afford professional accounting are partially excluded
from certain forms of economic activity (formal contracts,
audited transactions, disclosures that require an attestor).
A protocol that produces the accounting record as a byproduct
of ordinary settlement makes the record available to
participants who could not otherwise have afforded it. This
is a distributional effect: accounting infrastructure becomes
available at the protocol layer to participants who
previously faced the full cost of the apparatus. The
consequences --- for smaller economic units, for participants
in jurisdictions without strong accounting-profession
infrastructure, for machine and agent-led economic activity
whose scale does not justify an accountant --- are worth
working out in their own right, and we expect subsequent
work in the portfolio or by others to address them.

\section*{Acknowledgments}
I owe this paper's starting observation to an exchange with
a collaborator who framed the core contract, informally, as
``a balance sheet whose process is a P\&L.'' The framing
captured in plain language what would take a specialist
considerable technical vocabulary to articulate, and the
paper is substantially an elaboration of that framing into
the traditions of accounting research it engages. Pacioli
deserves most of the remaining credit; for the errors, I
alone.


% ════════════════════════════════════════════════════════════════════
% References
% ════════════════════════════════════════════════════════════════════

\begin{thebibliography}{99}

\bibitem[Braudel(1982)]{braudel1982wheels}
Braudel, F.
\newblock \emph{Civilization and Capitalism, 15th--18th Century,
  Volume 2: The Wheels of Commerce}.
\newblock Harper \& Row, New York, 1982.
\newblock Trans. S. Reynolds.

\bibitem[Buterin(2014)]{buterin2014ethereum}
Buterin, V.
\newblock Ethereum: A Next-Generation Smart Contract and
  Decentralized Application Platform.
\newblock White paper, 2014.
\newblock \url{https://ethereum.org/en/whitepaper/}

\bibitem[Chandler(1977)]{chandler1977visible}
Chandler, A.~D.
\newblock \emph{The Visible Hand: The Managerial Revolution
  in American Business}.
\newblock Belknap Press of Harvard University Press,
  Cambridge, MA, 1977.

\bibitem[Dai and Vasarhelyi(2017)]{dai2017toward}
Dai, J. and Vasarhelyi, M.~A.
\newblock Toward Blockchain-Based Accounting and Assurance.
\newblock \emph{Journal of Information Systems},
  31(3):5--21, 2017.

\bibitem[de Roover(1963)]{deroover1963medici}
de Roover, R.
\newblock \emph{The Rise and Decline of the Medici Bank,
  1397--1494}.
\newblock Harvard University Press, Cambridge, MA, 1963.

\bibitem[Grigg(2005)]{grigg2005triple}
Grigg, I.
\newblock Triple Entry Accounting.
\newblock Working paper, Systemics Inc., 2005.

\bibitem[Ijiri(1986)]{ijiri1986triple}
Ijiri, Y.
\newblock A Framework of Triple-Entry Bookkeeping.
\newblock \emph{Accounting, Organizations and Society},
  11(4--5):321--339, 1986.

\bibitem[Kindleberger(1993)]{kindleberger1993financial}
Kindleberger, C.~P.
\newblock \emph{A Financial History of Western Europe}.
\newblock Oxford University Press, New York, 2nd ed., 1993.

\bibitem[Nakamoto(2008)]{nakamoto2008bitcoin}
Nakamoto, S.
\newblock Bitcoin: A Peer-to-Peer Electronic Cash System.
\newblock White paper, 2008.
\newblock \url{https://bitcoin.org/bitcoin.pdf}

\bibitem[Pacioli(1494)]{pacioli1494summa}
Pacioli, L.
\newblock \emph{Summa de Arithmetica, Geometria,
  Proportioni et Proportionalita}.
\newblock Paganinus de Paganinis, Venice, 1494.
\newblock Particolaris de computis et scripturis (book on
  accounts and records), treatise XI.

\bibitem[Power(1997)]{power1997audit}
Power, M.
\newblock \emph{The Audit Society: Rituals of Verification}.
\newblock Oxford University Press, 1997.

\bibitem[Sombart(1916)]{sombart1916moderne}
Sombart, W.
\newblock \emph{Der moderne Kapitalismus}, volume 2.
\newblock Duncker \& Humblot, Munich, 1916.

\bibitem[Yamey(2005)]{yamey2005pacioli}
Yamey, B.~S.
\newblock Pacioli's \emph{Particularis de Computis et
  Scripturis}, 1494: Original Text and English Translation
  with Commentary.
\newblock In \emph{The Development of Accounting in the
  Late Medieval and Early Modern Periods}, Pacioli Society,
  2005.

\bibitem[Mattessich(1964)]{mattessich1964accounting}
Mattessich, R.
\newblock \emph{Accounting and Analytical Methods: Measurement and
  Projection of Income and Wealth in the Micro- and Macro-Economy}.
\newblock Richard D. Irwin, Homewood, IL, 1964.

\bibitem[Ijiri(1989)]{ijiri1989momentum}
Ijiri, Y.
\newblock \emph{Momentum Accounting and Triple-Entry Bookkeeping:
  Exploring the Dynamic Structure of Accounting Measurements}.
\newblock American Accounting Association, Sarasota, FL, 1989.

\bibitem[FASB(2014)]{fasbasc606}
Financial Accounting Standards Board.
\newblock \emph{Accounting Standards Codification Topic 606:
  Revenue from Contracts with Customers}.
\newblock FASB, Norwalk, CT, 2014.

\bibitem[IASB(2014)]{ifrs2014ifrs15}
International Accounting Standards Board.
\newblock \emph{IFRS 15: Revenue from Contracts with Customers}.
\newblock IFRS Foundation, London, 2014.

\bibitem[AICPA(2020a)]{aicpa2020auc315}
American Institute of Certified Public Accountants.
\newblock \emph{AU-C Section 315: Understanding the Entity and Its
  Environment and Assessing the Risks of Material Misstatement}.
\newblock AICPA, Durham, NC, 2020.

\bibitem[AICPA(2020b)]{aicpa2020auc500}
American Institute of Certified Public Accountants.
\newblock \emph{AU-C Section 500: Audit Evidence}.
\newblock AICPA, Durham, NC, 2020.

\bibitem[ERC-1271(2018)]{eip1271}
Giordano, F., Lockyer, M., and Castonguay, P.
\newblock EIP-1271: Standard Signature Validation Method for
  Contracts.
\newblock Ethereum Improvement Proposal 1271, July 2018.
\newblock \url{https://eips.ethereum.org/EIPS/eip-1271}

\bibitem[Safe(2024)]{safe_doc_2024}
Safe.
\newblock \emph{Safe Smart Account Documentation}.
\newblock Safe Foundation, ongoing as of 2024.

\bibitem[Compound(2020)]{compound_governor_2020}
Compound Labs.
\newblock \emph{Compound Governor: On-Chain Governance Module}.
\newblock Compound Protocol Documentation, 2020.

\bibitem[OpenZeppelin(2024)]{openzeppelin_governor_2024}
OpenZeppelin.
\newblock \emph{Governor Contracts: Reference Implementation for
  On-Chain Governance}.
\newblock OpenZeppelin Contracts Documentation, ongoing as of 2024.

\bibitem[Sablier(2023)]{sablier_2023}
Sablier Labs.
\newblock \emph{Sablier Protocol: Token Streaming Infrastructure}.
\newblock Sablier Documentation, 2023.

\bibitem[Superfluid(2023)]{superfluid_2023}
Superfluid Finance.
\newblock \emph{Superfluid Protocol: Continuous Money Streams}.
\newblock Superfluid Documentation, 2023.

\bibitem[Messari(2023)]{messari_dao_treasury_2023}
Messari Research.
\newblock \emph{State of DAO Treasuries 2023}.
\newblock Messari Quarterly Report, 2023.

\bibitem[DeepDAO(2024)]{deepdao_treasury_2024}
DeepDAO.
\newblock \emph{DAO Treasury Aggregate Statistics}.
\newblock DeepDAO Analytics Platform, ongoing as of 2024.

\bibitem[Yermack(2017)]{yermack2017corporate}
Yermack, D.
\newblock Corporate Governance and Blockchains.
\newblock \emph{Review of Finance}, 21(1):7--31, 2017.

\end{thebibliography}

\end{document}
